% Данный файл распространяется под лицензией CC BY 4.0. Текст лицензии размещён на https://creativecommons.org/licenses/by/4.0/
% (c) Кафедра системного программирования, 2025

\documentclass[a4paper]{article}

\usepackage[a4paper, top=8mm, bottom=8mm, left=8mm, right=8mm]{geometry}

\usepackage{polyglossia}
\setdefaultlanguage[babelshorthands=true]{russian}
\setotherlanguage{english}

\usepackage{fontspec}
\setmainfont{FreeSerif}
\newfontfamily{\russianfonttt}[Scale=0.7]{DejaVuSansMono}

\usepackage[tiny, compact]{titlesec}

\usepackage{titling}
\setlength{\droptitle}{-1cm}
\pretitle{\begin{center}\begin{bfseries}\Large}
\posttitle{\par\end{bfseries}\end{center}}
\preauthor{\begin{center}\normalsize}
\postauthor{\par\end{center}\vspace{-1.8cm}}

\usepackage{hyperref}
\usepackage{bookmark}
\usepackage{csquotes}
\usepackage{float}
\usepackage[useregional]{datetime2}
% Сюда можно дописывать нужные вам пакеты.

\title{UCFS: Доработка репозитория}

\author{Щелокаев Григорий Сергеевич}

% Дата много места занимает, её лучше не указывать.
\date{}

\begin{document}

\maketitle

\begin{flushright}
    Группа: \emph{\enquote{25.Б72-мм}}

    Кафедра: \emph{Системного программирования}

    % Степень/звание и должность научника мы лучше вас знаем, их можно не указывать.
    Научный руководитель: \emph{Григорьев Семён Вячеславович}

    Номер семестра практики: \emph{2}
\end{flushright} 


\section{Постановка задачи}

Работа выполняется в рамках проекта UCFS\footnote{Ссылка на UCFS: \url{https://github.com/FormalLanguageConstrainedPathQuerying/UCFS} (дата обращения: \DTMdate{2026-06-14}).}. UCFS~--- это инструмент для решения задач анализа графов с дополнительными ограничениями на пути, заданными контекстно-свободной грамматикой какого-либо языка. Целью работы является обеспечение проекта соответствующей инфраструктурой (CI, публикация пакетов, контроль тестового покрытия).

Результаты работы необходимы для снижения временных издержек разработчиков проекта на регулярное слияние кода, ручной запуск тестов, публикацию jar-файлов, а также минимизации рисков при помощи контроля тестового покрытия. Для конечных пользователей продукта это полезно снижением рисков обнаружения некорректной работы программы. В выполнении работы в первую очередь заинтересованы разработчики данного инструмента.
\section{Описание предлагаемого решения}


Далее описано решение каждой конкретной задачи: автоматизация запуска тестов, контроль тестового покрытия, публикация пакетов:
\begin{itemize}
    \item Для задачи~--- автоматический запуск тестов~--- основной идеей решения является использование существующих платформ непрерывной интеграции, для этого необходимо в конфигурационных файлах задать последовательность определённых команд, которые в конечном счёте и будут являться той самой автоматизацией запуска тестов. Также стоит задать необходимые условия, при которых данные команды будут запускаться: push/pull request.
    \item Основной идеей решения такой задачи, как контроль тестового покрытия, является грамотное использование существующих плагинов, предназначенных специально для многомодульных Java/Kotlin проектов, а также установление порогов тестового покрытия. Под грамотным использованием плагинов подразумевается их корректное применение в проекте, где тесты выделены в отдельный модуль. Чтобы получаемые результаты о покрытии соответствовали реальности, необходимо явно указывать файлы с исходным кодом как тестов, так и самого инструмента, файлы с результатами тестов и другие данные, необходимые для составления отчётов. Реализация идеи заключается в изменении уже существующих файлов сборки отдельных модулей проекта.
    \item Основная идея решения такой задачи, как публикация пакетов: вынести публикацию в отдельный модуль-агрегатор, который объединяет скомпилированный код рабочих модулей в один артефакт, а ссылки на внешние зависимости выносит в POM-файл.
\end{itemize}

Для выполнения поставленных задач соответствующие части данного проекта были написаны на языке программирования Kotlin Script и языке сериализации данных YAML.
Реализация решения осуществлялась с использованием таких программных продуктов, как платформа непрерывной интеграции GitHub Actions, система сборки проектов Gradle, инструмент контроля тестового покрытия JaCoCo (Java Code Coverage).

\section{Верификация}
На данный момент, решение поставленных задач находится на стадии рецензирования разработчиками инструмента. Замечания, рекомендации и комментарии размещены в данном pull request~--- \url{https://github.com/vkutuev/UCFS/pull/4} (дата обращения: \DTMdate{2026-06-15}). На данный момент единственным способом подтверждения того, что проделанная работа удовлетворяет изначальным требованиям, является самостоятельная верификация решения.

Проверка выполненных настроек CI/CD, а также контроля тестового покрытия и публикации пакетов проводилась локально и на GitHub Actions в форке~--- \url{https://github.com/GRIGAeo/UCFS} (дата обращения: \DTMdate{2026-06-14}).

Для проверки корректности контроля тестового покрытия из корня проекта были запущены следующие команды.
\begin{itemize}
    \item ./gradlew solver:jacocoTestReport generator:jacocoTestReport~--- данная команда составляет отчёты о тестовом покрытии для модулей solver и generator в формате .csv, .html и .xml. Если проект по какой-то причине не собирается или не проходятся тесты, то команда падает с ошибкой.
    \item ./gradlew solver:jacocoTestCoverageVerification generator:jacocoTestCoverageVerification~--- это команда так же, как и предыдущая составляет соответствующие отчёты, но помимо этого устанавливает пороги тестового покрытия (т.~е. если реальные проценты по покрытию ниже, чем установленные, сборка закончится с ошибкой). На данный момент оба модуля generator и solver не соответствуют заранее оговорённым с научным руководителем требованиям по порогу тестового покрытия, поэтому данный процесс падает с ошибкой. При занижении порогов тестового покрытия до минимальных значений программа, запущенная данной командой, корректно завершается.
\end{itemize}

Для проверки корректности настроенного CI/CD использовалась платформа GitHub Actions в данном форке~--- \url{https://github.com/GRIGAeo/UCFS/tree/test-infrastructure} (дата обращения: \DTMdate{2026-06-14}). Далее будут представлены ссылки на различные сценарии того, как работал workflow \foreignquote{english}{Run tests and test coverage}, в зависимости от того, в каком состоянии находится проект.

\begin{enumerate}
    \item Результат работы workflow, если не проходится какой-то из тестов~--- \url{https://github.com/GRIGAeo/UCFS/actions/runs/27426831720/job/81066687011} (дата обращения: \DTMdate{2026-06-15}). На шаге Run tests в строчках 49--54 видно, что тест TestDotParser не прошёл, поэтому весь workflow закончился с ошибкой.
    \item Результат работы workflow, если все тесты прошли, но тестовое покрытие ниже установленных порогов~--- \url{https://github.com/GRIGAeo/UCFS/actions/runs/27427297190/job/81068283493} (дата обращения: \DTMdate{2026-06-15}). На шаге Run tests в строчках 70--75 видно, что тестовое покрытие модуля solver ниже, чем установленные пороги, поэтому workflow завершился с ошибкой.
    \item Результат работы workflow, если все тесты прошли и тестовое покрытие не ниже установленных порогов~--- \url{https://github.com/GRIGAeo/UCFS/actions/runs/27429060075} (дата обращения: \DTMdate{2026-06-15}). Здесь же видно, что все тесты прошли, покрытие не ниже установленных порогов и была составлена таблица о тестовом покрытии.
\end{enumerate}

Для проверки корректности конечного jar-файла со всеми скомпилированными .class-файлами и POM-файла с указанием ссылок на зависимости проекта локально из корня проекта была запущена следующая команда.
\begin{itemize}
    \item ./gradlew :publisher:publishToMavenLocal~--- в результате получается jar-файл, содержащий все скомпилированные файлы исходного кода без внешних библиотек, и POM-файл с указанием ссылок на все 6 зависимостей проекта.
\end{itemize}

Для проверки корректности публикации пакетов в очередной раз использовалась платформа GitHub Actions в данном форке \url{https://github.com/GRIGAeo/UCFS} (дата обращения: \DTMdate{2026-06-14}). Далее будут представлены варианты завершения workflow \foreignquote{english}{Publish package to GitHub Packages}, который запускается только при выходе новых версий проекта.
\begin{enumerate}
    \item Результат работы workflow, если не проходится какой-то из тестов~--- \url{https://github.com/GRIGAeo/UCFS/actions/runs/27437439879/job/81102848039} (дата обращения: \DTMdate{2026-06-15}). На шаге Run tests в строчках 48--53 указано, что тест TestDotParser не прошёл, поэтому и весь workflow упал с ошибкой.
    \item Результат работы workflow, если прошли все тесты~--- \url{https://github.com/GRIGAeo/UCFS/actions/runs/27438411312} (дата обращения: \DTMdate{2026-06-15}). Также по данной ссылке \url{https://github.com/GRIGAeo/UCFS/packages/3074361} (дата обращения: \DTMdate{2026-06-15}) размещён получившийся package, удовлетворяющий выше перечисленным условиям к публикуемым файлам.
\end{enumerate}

Вывод: все компоненты CI/CD корректно настроены и проверены~--- тесты запускаются, покрытие контролируется, публикация в GitHub Packages работает.

\section{Заключение}

В ходе выполнения работы получены следующие результаты:

\begin{itemize}
    \item После обсуждения с одним из разработчиков инструмента и по совместительству научным руководителем были установлены следующие пороги тестового покрытия: INSTRUCTION~--- 0.95, BRANCH~--- 0.8, LINE~--- 0.8, METHOD~--- 0.85, CLASS~--- 0.9.
    \item Результаты тестового покрытия для модуля solver приведены в таблице~\ref{tab:solver-coverage}.
    \item Для модуля generator тестов нет, поэтому все показатели равны 0.
    \item Для публикации jar-файлов был создан и корректно настроен отдельный модуль publisher, который публикует в GitHub Packages скомпилированные файлы проекта, а в POM-файле указаны все ссылки на внешние зависимости проекта.
    \item В файлах .github/workflows/ci-publishing-jar-files.yaml, .github/workflows/ci-test-infrastructure.yaml корректно реализованы процессы CI/CD с запуском тестов, проверкой порогов тестового покрытия и публикацией пакетов. 
\end{itemize}

Материалы, иллюстрирующие выполнение работы, отсутствуют.

Ссылка на pull request: \url{https://github.com/vkutuev/UCFS/pull/4} (дата обращения: \DTMdate{2026-06-14}).

Ссылка на fork: \url{https://github.com/GRIGAeo/UCFS} (дата обращения: \DTMdate{2026-06-14}).

Решение находится на рецензировании у разработчиков. После получения комментариев о проделанной работе и добавления последних корректировок изменения будут добавлены в основной репозиторий UCFS.

\begin{table}[H]
\centering
\begin{tabular}{|l|c|c|}
\hline
Показатель & Покрытие & Порог \\
\hline
INSTRUCTION & 58\% & 95\% \\
BRANCH      & 43\% & 80\% \\
LINE        & 54\% & 80\% \\
METHOD      & 53\% & 85\% \\
CLASS       & 70\% & 90\% \\
\hline
\end{tabular}
\caption{Тестовое покрытие модуля solver}
\label{tab:solver-coverage}
\end{table}

\end{document}
